\chapter{Espacios normados}

% En primer lugar recordaremos algunos conceptos de notación

\section{Espacio vectorial}

\begin{notacion}
    A lo largo del curso los espacios vectoriales serán reales o complejos y notaremos indistintamente por $\bb{K}$ al cuerpo.\\

    Sea $X$ un espacio vectorial y sean $A,B\subseteq X$ notaremos por $A+B$ al siguiente conjunto
    \begin{gather*}
        A+B = \{a + b \ : \ a\in A,\ b\in B\}
    \end{gather*}
    En el caso particular de que $A$ sea unitario, es decir, $A=\{x\}$, entonces notaremos 
    \begin{gather*}
        A + B = x + B
    \end{gather*} 
    en lugar de $\{x\} + B$.\\

    Dado $K\subseteq \bb{K}$ un subconjunto notaremos
    \begin{gather*}
        KA = \{\alpha a_i \ : \ \alpha \in K,\ a\in A\} \subseteq X
    \end{gather*}
    Análogamente al caso anterior de que $K$ sea unitario con $K=\{\alpha\}$ notaremos 
    \begin{gather*}
        KA = \alpha A
    \end{gather*}
\end{notacion}

\begin{definicion}
    Dado un subconjunto $M\subseteq X$ diremos que es un subespacio vectorial de $X$ si es cerrado para la suma y el producto, es decir
    \begin{gather*}
        M + M \subseteq M\\
        \bb{K} M \subseteq M
    \end{gather*}
\end{definicion}

\begin{definicion}
    Si $\emptyset \neq A \subseteq X$, el \textbf{subsepacio vectorial} de $X$ generado por $A$ lo denotamos por $Lin(A)$, que será el \textbf{envolvente lineal} de $A$ y vendrá dado por
    \begin{gather*}
        Lin(A) = \{\alpha_1 a_1 + ... + \alpha_na_n : n\in \bb{N},\ \alpha_1,...,\alpha_n\in \bb{K},\ a_1,...,a_n\in A\}
    \end{gather*}
\end{definicion}

\begin{definicion}
    Diremos que $A$ es \textbf{linealmente independiente} si ningún elemento de $A$ se puede poner como combinación lineal del resto.
\end{definicion}

\begin{definicion}
    Diremos que $A$ es una \textbf{base} de $X$ si es linealmente independiente y su envolvente lineal es $X$, es decir, $Lin(A)=X$.
\end{definicion}

\begin{definicion}
    Todo espacio vectorial $X$ tiene una base\footnote{Hablamos de bases algebraicas o de Hamel. Nosotros trabajaremos con bases de Hilbert}. Además, todas las bases de $X$ tienen el mismo cardinal. Este cardinal se llama \textbf{dimensión} del espacio $X$.
\end{definicion}

\section{Espacios normados}

\begin{definicion}
    Sea $X$ un espacio vectorial. Una \textbf{seminorma} sobre $X$ es una función $\varphi: X \to \bb{R}$ que verifica las siguientes propiedades
    \begin{enumerate}
        \item \underline{Desigualdad triangular:} $\varphi(x+y) \leq \varphi(x) + \varphi(y)$ para todo $x,y\in X$.
        \item $\varphi(\lambda x) = |\lambda|\varphi(x)$ para todo $\lambda \in \bb{K}$, $\forall x \in X$.
    \end{enumerate}
\end{definicion}

\begin{observacion}
    La primera consecuencia que podemos concluir es 
    \begin{gather*}
        \varphi(0) = \varphi(0 \cdot 0) = 0 \cdot \varphi(0) = 0
    \end{gather*}
    La segunda consecuencia es, para cualquier $x\in X$
    \begin{gather*}
        0 =\varphi(0) = \varphi(x + (-x)) \leq \varphi(x) + |-1|\varphi(x) = 2\varphi(x) \Rightarrow \varphi(x) \geq 0
    \end{gather*}
\end{observacion}

\begin{definicion}
    Una \textbf{norma} es una seminorma que verifica además
    \begin{gather*}
        \|x\| = 0 \sii x=0
    \end{gather*}
\end{definicion}

\begin{definicion}
    Al par $(X, \|\cdot\|)$ donde $X$ es un espacio vectorial y $\|\cdot\|$ es una norma se le llama \textbf{espacio normado}.
\end{definicion}

\begin{notacion}
    Por lo general nos referiremos a un espacio normado $(X, \|\cdot\|)$ por $X$ e implícitamente tendrá asociado una norma.
\end{notacion}

\begin{definicion}
    Sea $X$ un espacio normado, $a\in X$ y sea $r\in \bb{R}^+$. Definimos la \textbf{bola abierta} de centro $a$ y radio $r$ como 
    \begin{gather*}
        B(a,r) = \{x\in X \ : \ \|a-x\| < r\}
    \end{gather*}
    Definimos la \textbf{bola cerrada} de centro $a$ y radio $r$ como 
    \begin{gather*}
        \bar{B}(a,r)= \{x\in X \ : \ \|a-x\| < r\}
    \end{gather*}
    Definimos la \textbf{esfera} de centro $a$ y radio $r$ como 
    \begin{gather*}
        S(a,r) = \{x\in X \ : \ \|a-x\| < r\}
    \end{gather*}
\end{definicion}

\begin{definicion}
    Un subconjunto $A\subset X$ se dice \textbf{abierto} si 
    \begin{gather*}
        \forall a \in A \ \ \exists r \in \bb{R}^+ : B(a,r)\subseteq A
    \end{gather*}
\end{definicion}

\begin{observacion}
    Esta definición es equivalente a considerar que $A$ es abierto si es unión arbitraria de bolas abiertas.
\end{observacion}

\begin{definicion}
    Esta familia de conjuntos es una topología en $X$ a la que llamaremos \textbf{topología de la norma}. Además, con esta definición tendremos que $\forall a \in X$, el conjunto $\{B(a,r) \ : \ r>0\}$ es una base de entornos de $a$.
\end{definicion}

\begin{notacion}
    Notaremos $B = \{x\in X \ : \ \|x\| < 1\}$ y de esta forma cualquier bola la podremos escribir como 
    \begin{gather*}
        B(a,r) = a + rB
    \end{gather*}
\end{notacion}

\begin{definicion}
    Sea $X$ un espacio normado y $A\subset X$. Se dice que $A$ está \textbf{acotado} si existe un $M\in \bb{R}$ tal que 
    \begin{gather*}
        \|x\| < M \quad \forall x \in A
    \end{gather*}
\end{definicion}